% Generated draft: requires mathematical and source review.
\documentclass[11pt]{article}
\usepackage[margin=25mm]{geometry}
\usepackage{fontspec}
\setmainfont{Latin Modern Roman}
\usepackage{amsmath,amssymb,mathtools,unicode-math}
\setmathfont{Latin Modern Math}
\usepackage{enumitem,needspace,xurl,hyperref}
\setlength{\parindent}{0pt}
\setlength{\parskip}{5pt}
\setlength{\emergencystretch}{3em}
\newcommand{\N}{\mathbb{N}}
\newcommand{\Z}{\mathbb{Z}}
\newcommand{\Q}{\mathbb{Q}}
\newcommand{\R}{\mathbb{R}}
\newcommand{\C}{\mathbb{C}}
\newcommand{\F}{\mathbb{F}}
\newcommand{\A}{\mathbb{A}}
\newcommand{\PP}{\mathbb{P}}
\newcommand{\E}{\mathbb{E}}
\newcommand{\eps}{\varepsilon}
\newcommand{\dd}{\,\mathrm{d}}
\begin{document}
\section*{UnsolvedMath ID: 30000306}
\textbf{Source problem code:} catalogue entry
\subsection{Definitions and mathematical statement}
Let $\pi_\theta$ be the Dirichlet or Poisson--Dirichlet stationary probability in the model of the report, on its specified state space $S_\theta$, and let $Q_n^{(\theta)}(x,y)$ be the corresponding expansion kernels. The restored positivity question concerns
\[
 P(x,\dd y)=\left[1+\sum_{n\ge1}a_n(\theta)
                              Q_n^{(\theta)}(x,y)\right]\pi_\theta(\dd y).
\]
It asks to characterize the nonzero coefficient sequences for which this defines a nonnegative kernel. In the usual centered normalization, $\int Q_n^{(\theta)}(x,y)\,\pi_\theta(\dd y)=0$, so the bracket integrates to one when summation and integration are justified. The convergence and normalization conventions for $Q_n$ must be supplied; positivity of a merely formal series is not enough to define a transition kernel.

The second question asks whether the standard diffusion coefficients $\rho_n^{(|\theta|)}(t)$ are the only time-dependent choices $a_n(t,\theta)$ yielding a Markov transition function with that stationary measure. Beyond positivity at each $t>0$, this requires
\[
 P_{t+s}(x,B)=\int_{S_\theta}P_t(x,\dd z)P_s(z,B),\qquad
 \int P_t(x,B)\,\pi_\theta(\dd x)=\pi_\theta(B),
\]
for all measurable $B$, together with the intended identity-at-zero and continuity conventions. Under an orthogonal-projection-kernel normalization, the semigroup condition is expressed by multiplicativity of the spectral coefficients, with the requisite convergence assumptions; the actual kernel convention is not specified in the extract.

The Dirichlet parameter domain and finite state simplex, the one- or multiparameter Poisson--Dirichlet model, $Q_n$ normalization, meaning of ``nonzero'', and formula for $\rho_n$ are not given. Neither is a fixed generator or time normalization. This last omission matters for uniqueness: replacing any semigroup $P_t$ by $P_{ct}$ for $c>0$ preserves the Markov and stationary properties while changing its time-dependent coefficients. Thus the literal uniqueness target needs restrictions or an equivalence convention; no unsupported universal uniqueness claim is made.

The literature assessment checked on 21 August 2026 attributes complete positivity criteria in important Jacobi and Gegenbauer parameter ranges to Bochner and Gasper, while saying the all-parameter request is broader. The precise special ranges and theorem statements are absent. This is supplied dated partial-progress triage requiring independent verification.
\subsection{Short English statement}
Classify coefficients making the Dirichlet or Poisson--Dirichlet kernel expansion positive, and then impose the semigroup conditions needed for a Markov process. Uniqueness of the standard diffusion coefficients requires a generator or time-normalization restriction that the extract does not specify.
\subsection{Sources}
\begin{itemize}
\item \emph{Mathematical Population Genetics} (2006), Oberwolfach Reports, DOI 10.4171/owr/2005/40, pp. 2241--2304, \url{https://doi.org/10.4171/owr/2005/40}; proposed year supplied as 2005.
\item The assessment checked 21 August 2026 cites only the original report, \url{https://doi.org/10.4171/owr/2005/40}, while naming Bochner and Gasper's special-range positivity results without full additional references.
\item Local UnsolvedMath registry, permanent ID $30000306$, exact source code \texttt{catalogue entry}, created 25 August 2026 and updated 16 September 2026. Supplied verification restores the kernel expansion, coefficient sequences and uniqueness target.
\end{itemize}
\end{document}
